\section{An intrinsic detector theorem on a compact double}
\label{sec:v86-meromorphic}
The polar tests have a formulation which does not require a rational
coordinate. The ambient object is a compact Riemann surface with a real
structure. Periods enter the obstruction test, and the canonical divisor
enters the clock budget. Thus principal parts alone no longer give the
answer.

Let $X$ be a connected compact Riemann surface of genus $g$, equipped
with an antiholomorphic involution $\sigma$. Choose disjoint open arcs
$I,J$ in its real locus and a real point $b\notin I\cup J$. A real
analytic parameter on $I$ is fixed for differentiating action locations.
The arc $J$ is the clock domain; choose a base clock $T_*\in J$.
A differential $\omega$ is called real if
$\sigma^*\overline\omega=\omega$. An imaginary-period normalization
means that the real part of every period on the punctured surface is
zero. The terminology specifies the normalization, not the location of
its residues: the residues below are real.

\begin{lemma}[Normalized action differentials]\label{lem:v86-normalized}
For distinct real points $x,b$, there is a unique meromorphic
differential $\omega_{x,b}$ with simple poles of residues $+1,-1$
at $x,b$, no other poles, and imaginary periods. It is real and depends
smoothly on $x$ away from $b$. The differential
\[
                         B_x=\partial_x\omega_{x,b}
\]
has a nonzero double principal part at $x$, zero residue, no other
poles, and imaginary periods. It is independent of the choice of $b$.
A change of local action parameter rescales $B_x$ by a nonzero real
factor.
\end{lemma}
\begin{proof}
The prescribed-principal-part theorem for differentials supplies a
meromorphic differential with these two residues: their sum is zero.
This is the usual Riemann--Roch consequence, recalled in
\cite{V86Forster}. Choose a homology basis of $2g$ cycles avoiding the
poles. The real-linear map from holomorphic differentials to the real
parts of their periods on these cycles is injective. Indeed, if these
real parts vanish, the real part of a local primitive is a globally
defined harmonic function on $X$, hence constant. Its differential is
zero, so the holomorphic differential is zero. Both real dimensions
are $2g$, and the map is therefore an isomorphism. Subtract the unique
holomorphic differential cancelling the real parts of the prescribed
periods. Small loops around the poles already have imaginary periods,
because their integrals are $2\pi i$ times real residues. This proves
existence. The same harmonic argument proves uniqueness.

Conjugation by $\sigma$ preserves the specified poles, residues and
normalization, and hence preserves the differential by uniqueness.
The local principal-part construction and the fixed finite-dimensional
period correction depend smoothly on $x$ in a coordinate neighbourhood;
uniqueness joins these constructions. In a local coordinate $z$ the
moving simple principal part is $\dd z/(z-x)$. Its derivative has
principal part $\dd z/(z-x)^2$. The principal part at $b$ is fixed,
so differentiation removes it. All periods remain imaginary. The
difference obtained from two choices of $b$ is a holomorphic
differential with imaginary periods and must vanish. The chain rule
proves the final assertion.
\end{proof}
The imaginary-period convention is the real-normalized convention
applied after multiplication by $i$. Normalized differentials on
compact and degenerating curves are studied, for example, in
\cite{V86GKN}; only the fixed-surface existence and uniqueness just
proved are used here.

Define the action curve on the clock arc by
\begin{equation}\label{eq:v86-intrinsic-curve}
            L_x(T)=\int_{T_*}^{T}\omega_{x,b},\qquad T\in J.
\end{equation}
The integral is along $J$ and is real. Let $D$ be an effective divisor
whose support avoids $J$. Let $E$ be a finite-dimensional real subspace
of the real meromorphic differentials with poles bounded by $D$, and
put
\begin{equation}\label{eq:v86-intrinsic-detectors}
 \mathcal F_E=\left\{c+\int_{T_*}^{T}e:\ c\in\R,
                                             \ e\in E\right\}.
\end{equation}
Constants are included, and the primitive is only required on $J$.
The observed scalar law is $f=e+L_x-L_y$, with $e\in\mathcal F_E$
and ordered distinct $x,y\in I$. We call this detector space regularly
separating if this representation is unique and its parameter map into
functions on $J$ has injective differential at every such pair. In
particular this definition includes the nuisance coefficients.

Let $S=\operatorname{supp}D\cap I$ and define the finite charge set
\[
 \mathcal Z_2(S)=\{c\in\Z^S\setminus\{0\}:\ \sum_s c_s=0,
                                      \ \sum_s\max(c_s,0)\le2\}.
\]
For $c\in\mathcal Z_2(S)$ let $\omega_c$ be the unique
imaginary-period differential with simple principal parts of residues
$c_s$ at $s\in S$, and no other poles. Equivalently,
$\omega_c=\sum_s c_s\omega_{s,b}$; the pole at $b$ cancels.

\begin{theorem}[Divisor--period classification]\label{thm:v86-divisor}
The detector space $\mathcal F_E$ is regularly separating if and only if
\begin{align}
 \omega_c&\notin E &&(c\in\mathcal Z_2(S)),
                                      \label{eq:v86-charge-test}\\
 E\cap\operatorname{span}_{\R}\{B_s:s\in A\}&=\{0\}
             &&(A\subset S,\ 1\le |A|\le2).
                                      \label{eq:v86-double-test}
\end{align}
These are finitely many membership and intersection tests in a
finite-dimensional space of meromorphic differentials. They use the
full normalized differentials, including their holomorphic period
corrections, rather than just their principal parts.
\end{theorem}
\begin{proof}
A scalar collision, differentiated on the clock arc, gives a member of
$E$ equal, up to sign, to
\[
       \omega_{x,b}-\omega_{y,b}-\omega_{x',b}+\omega_{y',b}.
\]
By analytic continuation the equality is one of meromorphic
differentials on $X$. Its residue divisor is
$\delta_x-\delta_y-\delta_{x'}+\delta_{y'}$.
Every noncancelled pole must lie in $S$. If this divisor is nonzero,
its coefficients belong to $\mathcal Z_2(S)$, and normalization
identifies the differential with $\omega_c$. If the divisor is zero,
equality of the signed measures $\delta_x-\delta_y$ and
$\delta_{x'}-\delta_{y'}$, with both pairs off the diagonal, gives
$x=x'$ and $y=y'$. The remaining nuisance difference is constant,
and the original collision makes that constant zero.

Conversely every $c\in\mathcal Z_2(S)$ is such a secant divisor.
For positive mass two, list its positive atoms with multiplicity as
$p_1,p_2$ and its negative atoms as $n_1,n_2$; use the ordered pairs
$(p_1,n_1)$ and $(n_2,p_2)$. For positive mass one, write
$c=\delta_p-\delta_n$ and use $(p,z),(n,z)$, where
$z\in I\setminus\{p,n\}$. Both pairs are nondegenerate.
If $\omega_c\in E$, integration and the included constants give a
scalar collision. This proves the exact-identification criterion.

A parameter tangent, differentiated in the clock variable, is a member
of $E$ plus $uB_x-vB_y$. Nonzero double principal parts at distinct
points cannot cancel. Every action with nonzero tangent coefficient
must therefore belong to $S$. The tangent obstruction is precisely
\eqref{eq:v86-double-test}; a one-point obstruction is realized by
choosing any different second action with zero tangent coefficient.
Once again constants remove the integration ambiguity. If no
obstruction occurs, both action coefficients vanish, and the original
zero tangent then forces every nuisance coefficient to vanish.
\end{proof}
The theorem applies on positive-genus surfaces even though no global
rational coordinate exists. For example, take $E$ to be any real
subspace of holomorphic differentials on $X$ and $D=0$. Both tests
are empty, so all these detector spaces are regularly separating.
More generally, differentials with poles only outside the action arc
are admissible. Coupled subspaces with poles on $I$ require the full
tests: replacing a normalized differential by another differential
with the same principal parts can change its membership in $E$.

\begin{theorem}[A divisor clock budget]\label{thm:v86-divisor-budget}
If the tests of Theorem~\ref{thm:v86-divisor} hold, any
\begin{equation}\label{eq:v86-divisor-budget}
                            N=\deg D+2g+4
\end{equation}
distinct clocks in $J$ identify the nuisance and ordered unequal
actions, and the finite evaluation map is an immersion. Its inverse
is locally smooth. On every compact set of coefficient and action
parameters away from the diagonal and the ends of the coordinate
arcs, it has a uniform Lipschitz constant.
\end{theorem}
\begin{proof}
Let $H$ be a difference of two scalar laws, or a scalar parameter
tangent. If it vanishes at $N$ ordered clocks, Rolle's theorem on the
real arc gives at least $N-1$ distinct zeros of its differential.
This differential extends meromorphically to $X$. Its poles are bounded
by $D$ together with at most four additional pole orders: four simple
moving poles for a secant, or two double moving poles for a tangent.
The auxiliary pole at $b$ cancels. Its pole degree is at most
$\deg D+4$.

A nonzero meromorphic differential on a compact surface of genus $g$
has zero degree minus pole degree equal to $2g-2$. Thus the number of
its distinct zeros is at most $\deg D+2g+2$. Since
$N-1=\deg D+2g+3$, the differential must vanish identically.
A sampled zero then gives $H=0$. The two tests in
Theorem~\ref{thm:v86-divisor} give injectivity and full differential
rank respectively. A nonsingular coordinate minor gives a smooth
local inverse. For completeness, on a compact parameter set the local
inverse estimates have a finite subcover. Pairs not covered by a
common sufficiently small local neighbourhood have a positive minimum
observed separation by compactness and injectivity. Combining the two
bounds gives the asserted global Lipschitz estimate on that set.
\end{proof}
This is a sufficient budget, not a universal minimal-clock theorem.
If $d=\dim\mathcal F_E$, immersion necessarily requires $N\ge d+2$.
There can be a gap between this dimensional lower bound and
\eqref{eq:v86-divisor-budget}. A smaller optimal count requires extra
structure; the polynomial oriented-interval theorem supplies exactly
such structure. The two-clock shared-nuisance theorem in
Section~\ref{sec:v86-shared} instead changes the observation domain
and cannot be obtained by optimizing a pointwise zero count.

\subsection{The rational theorem as the genus-zero case}
On the sphere with its real structure,
$\omega_{x,b}=\dd T/(T-x)-\dd T/(T-b)$, and hence
$L_x-L_y$ is $\log(T-x)-\log(T-y)$ up to a constant on a clock
interval lying to the right of the actions. Also
$B_s=\dd T/(T-s)^2$. Theorem~\ref{thm:v86-divisor} therefore gives
the charge and double-pole classification for rational derivatives.

More precisely, let $D_f$ be the degree of the finite common
denominator of a rational derivative space, and let $\ell$ be its
largest order at infinity, with $\ell=-\infty$ for the zero space.
Passing from rational functions to differentials adds the possible
pole at infinity of order $\max\{\ell+2,0\}$. The minimal ambient
pole divisor has degree
\[
              \deg D=D_f+\max\{\ell+2,0\}.
\]
Thus \eqref{eq:v86-divisor-budget} becomes
$D_f+6+\max\{\ell,-2\}$, exactly the rational pole budget. The
additional $2g$ in the intrinsic theorem is the canonical-divisor
contribution, while the normalization in the obstruction test records
the period geometry. Neither contribution is visible in a list of
rational partial fractions.

\subsection{Predecessors and the additional statement}
Riemann--Roch, the residue constraint, period normalization and the
degree of the canonical divisor are classical; see
\cite{V86Forster,V86GKN}. The theorem does not propose a new existence
theorem for normalized differentials. Its additional statement is the
necessary-and-sufficient finite obstruction classification of unknown
detector subspaces for this action quotient, together with a finite
clock immersion theorem expressed in intrinsic divisor data. This
extends the rational criterion to nonrational compact doubles and
shows precisely where principal-part tests need period information.
